% Editable transcription; verify reused text against the original PDF.
\section{Discussion}

AI has seen multiple cycles of enthusiasm and disappointment, but
the current wave seems to be of a different order.
As we stated in the introduction of this paper, one of the original
goals of AI was to learn more about human intelligence.
This endeavor could be misguided as AI may only produce
``cognitive wheels'', techniques that have no equivalent in
human cognition.
In this paper we argued that this might have been true for some
older approaches (e.g. brute force search techniques), but is less
the case for much of current AI.
The progress made in recent years is certainly technologically
driven, but inspired by biological and psychological knowledge
about human information processing and learning.

We expect that the recent progress in AI will change the way we
think about intelligence.
AI forces us to rethink the definition of intelligence.
Definitions that center on just information processing and problem
solving are perhaps insufficient.
Shank's observation that intelligence is all about generalization
has, so far, withstood the test of time.
Many information processing problems, from processing speech
to playing chess, appear to be less difficult than perhaps expected.
The really hard problem is to deal with completely novel cases.
One requirement for solving this hard problem is the ability to
learn invariant and thus generalizable patterns.
And especially with regard to learning, the progress in AI has
been spectacular.
The main difference between AI systems of the past, such as expert
systems, and modern AI is the fact that they learn.
That deep learning and reinforcement learning, the core techniques
in current AI, have deep roots in psychology is remarkable and
promising for studying how artificial and human intelligence
are related.

AI is relevant to intelligence research because it enhances our
understanding of the core mechanisms of human cognition.
How the immense neural systems in our brain are able to process
extremely complicated information such as speech and produce
logical thinking is an extremely difficult question.
Having an artificial system that performs such tasks using the
same basic principles is extremely useful.
Classic questions regarding the modularity of the mind, the origin
of creativity, and the organization of long-term memory spring
to mind.
In addition, we argued that the psychological relevance of AI
extends to unexpected areas such as the understanding of
individual differences and the development of cognition.
It is relatively easy to create a population of AI systems with
minor differences in architecture and training regime.
Modern AI provides us with a new playing field for individual
differences research.

On a practical level we expect fruitful interactions regarding
the measurement of natural and artificial intelligence.
As modern AI systems are incredibly complex, our experience
in examining such systems may be relevant for AI.
Vice versa, insights from AI may lead to new developments in
(adaptive) intelligence testing and educational interventions.

We attempted to shed light on the future of intelligence research
from the point of view of AI.
Our overview is necessarily limited and probably quickly outdated,
but hopefully we have given intelligence researchers some insights
in the rapid developments in AI and the possible consequences
for our field.
